\section*{Hoofdstuk \ref{ch:b2:batteries-electrolysis} --- Batterijen, brandstofcellen en elektrolyse}

\begin{solution}{exo:b2:batteries-electrolysis:1}
Zink alleen: ongeveer \qty{-0.82}{V} en een stroom van 0.15 (modeleenheden). In contact
met koper: ongeveer \qty{-0.74}{V} en 2.4. De stroom, en dus de oplossnelheid, is
in contact met koper zo'n zestien keer groter.
\end{solution}

\begin{solution}{exo:b2:batteries-electrolysis:2}
$m = 500 \times 3600 \times 65.38/(2 \times 96\,485) = \qty{610}{g}$.
\end{solution}

\begin{solution}{exo:b2:batteries-electrolysis:3}
$E^\circ = \qty{2.04}{V}$ (zoals in het hoofdstuk); zes cellen in serie geven
\qty{12}{V}.
\end{solution}

\begin{solution}{exo:b2:batteries-electrolysis:4}
$1/6.94 = \qty{0.144}{mol}$; $0.144 \times 96\,485 = \qty{1.39e4}{C}$, dus
\qty{3.86}{A.h}.
\end{solution}

\begin{solution}{exo:b2:batteries-electrolysis:5}
Bij 100~\%: $200 \times 86\,400 \times 63.546/(2 \times 96\,485) = \qty{5.69}{kg}$;
$\eta_F = 5.50/5.69 = 0.967$.
\end{solution}

\begin{solution}{exo:b2:batteries-electrolysis:6}
$w = 2 \times 96\,485 \times 3.3/(0.90 \times 0.06538) = \qty{1.08e7}{J/kg}$, dus
\qty{3.0}{kW.h/kg}.
\end{solution}

\begin{solution}{exo:b2:batteries-electrolysis:7}
$U = E - rI$: $r = (1.55 - 1.40)/(0.60 - 0.10) = \qty{0.30}{\ohm}$; $E = 1.55 + 0.30
\times 0.10 = \qty{1.58}{V}$.
\end{solution}

\begin{solution}{exo:b2:batteries-electrolysis:8}
Anode \ce{2Cl- -> Cl2 + 2e-}; kathode \ce{2H2O + 2e- -> H2 + 2OH-}. $\Delta_r G^\circ =
2(-157.24) - 2(-131.23) - 2(-237.13) = \qty{422.2}{kJ/mol}$, $U_{\min} =
\qty{2.19}{V}$. Chloor zou met het hydroxide reageren (\ce{Cl2 + 2OH- -> Cl- +
ClO- + H2O}), wat beide producten bederft, en een mengsel van chloor en waterstof
kan ontploffen.
\end{solution}

\begin{solution}{exo:b2:batteries-electrolysis:9}
\ce{NaCl -> Na + 1/2Cl2}, $n = 1$: $U_{\min} = 310\,674/96\,485 = \qty{3.22}{V}$.
Per kilogram natrium: $96\,485 \times 3.22/0.02299 = \qty{1.35e7}{J}$, dus
\qty{3.75}{kW.h}.
\end{solution}

\begin{solution}{exo:b2:batteries-electrolysis:10}
Op een zinkkathode is de reductie van \ce{H+} traag en ligt haar muur onder de
zinkgolf: naarmate de potentiaal verlaagd wordt, zet zink zich eerst af, met weinig
waterstof. Op platina is de reductie van \ce{H+} snel en begint ze rond
\qty{0}{V}: in het begin ontwikkelt zich alleen waterstof. Zodra het platina met
zink bedekt is, gedraagt de kathode zich als een zinkkathode en zet zink zich af.
\end{solution}

\begin{solution}{exo:b2:batteries-electrolysis:11}
Bij \qty{1250}{K}: $\Delta_f G^\circ(\ce{Al2O3}) = -1328.3 + 0.75 \times 66.0 =
\qty{-1278.8}{kJ/mol}$, $\Delta_f G^\circ(\ce{CO2}) = \qty{-396.1}{kJ/mol}$. $\Delta_r G^\circ =
3(-396.1) - 2(-1278.8) = \qty{1369}{kJ}$, $n = 12$: $U_{\min} = \qty{1.18}{V}$. Inerte anode:
$\Delta_r G^\circ = \qty{2557.5}{kJ}$, $U_{\min} = \qty{2.21}{V}$; de verbrandende koolstof
levert ongeveer \qty{1.03}{V} van de arbeid.
\end{solution}

\begin{solution}{exo:b2:batteries-electrolysis:12}
Per mol water is de lading $2F$ in beide richtingen dezelfde: rendement $0.70/1.80 =
0.39$. De rest is warmte: de \omterm{def:b2:current-potential-curves:overpotential}{overspanningen} van beide zuurstofelektroden (traag in
beide richtingen), die van de waterstofelektroden, en de ohmse verliezen.
\end{solution}

\begin{solution}{pb:b2:batteries-electrolysis:1}
\textbf{1.} Aluminium $+3 \to 0$; zuurstof blijft $-2$; koolstof $0 \to +4$.
\textbf{2.} Kathode \ce{Al^3+ + 3e- -> Al}; anode \ce{C + 2O^2- -> CO2 + 4e-}.
\textbf{3.} \ce{2Al2O3 + 3C -> 4Al + 3CO2}, $n = 12$.
\textbf{4.} \ce{Al2O3}: $-1328.3 + 0.75 \times (1328.3 - 1262.3) = \qty{-1278.8}{kJ/mol}$;
\ce{CO2}: \qty{-396.1}{kJ/mol}.
\textbf{5.} $3(-396.1) - 2(-1278.8) = \qty{+1369}{kJ}$.
\textbf{6.} $1\,369\,000/(12 \times 96\,485) = \qty{1.18}{V}$.
\textbf{7.} $2 \times 1278.8/(12 \times 96.485) = \qty{2.21}{V}$: de oxidatie van koolstof,
die zelf \omterm{def:b2:reaction-free-energy:exergonic}{exergonisch} is, betaalt bijna de helft van de arbeid om het oxide te ontleden.
\textbf{8.} $10^6/26.982 = \qty{3.706e4}{mol}$.
\textbf{9.} $3 \times 96\,485 \times 3.706 \times 10^4 = \qty{1.073e10}{C}$.
\textbf{10.} $1.073 \times 10^{10}/0.94 = \qty{1.141e10}{C}$.
\textbf{11.} $1.141 \times 10^{10}/3.0 \times 10^5 = \qty{3.80e4}{s}$, \qty{10.6}{h}.
\textbf{12.} $24/10.6 = 2.27$: \qty{2.27}{t} per dag.
\textbf{13.} Een deel van het in het bad opgeloste aluminium bereikt het anodegas en
wordt door koolstofdioxide opnieuw geoxideerd, zodat zijn lading tweemaal besteed wordt; er
lekt ook wat stroom door het bad zonder het te elektrolyseren.
\textbf{14.} $1.141 \times 10^{10} \times 4.1 = \qty{4.68e10}{J}$, \qty{13.0}{MW.h}.
\textbf{15.} \qty{13.0}{kW.h/kg}.
\textbf{16.} $3 \times 96\,485 \times 1.18/0.026982 = \qty{1.27e7}{J/kg}$, \qty{3.5}{kW.h/kg}.
\textbf{17.} Naar warmte: de \omterm{def:b2:current-potential-curves:overpotential}{overspanningen} van de elektroden en vooral de
ohmse val in het bad. Die warmte houdt het bad gesmolten bij ongeveer
\qty{1250}{K}, waarvoor anders brandstof nodig zou zijn.
\textbf{18.} $4.1 \times 3.0 \times 10^5 = \qty{1.23}{MW}$.
\textbf{19.} $28\,690/0.026982 = \qty{1.06e6}{J/kg}$, \qty{0.30}{kW.h/kg}: ongeveer 2~\% van
de energie van de \omterm{def:b2:batteries-electrolysis:electrolysis}{elektrolyse}.
\textbf{20.} $1.18/4.1 = 0.29$.
\textbf{21.} $3.706 \times 10^4 \times 3/4 \times 12.011 = \qty{334}{kg}$.
\textbf{22.} $3.706 \times 10^4 \times 3/4 \times 44.009 = \qty{1.22e3}{kg}$.
\textbf{23.} De hete bovenkanten van de anoden verbranden in de lucht; een deel van de koolstof vertrekt
als koolstofmonoxide (\ce{C + CO2 -> 2CO} bij die temperatuur), wat per zuurstofatoom tweemaal
zoveel koolstof kost.
\textbf{24.} \qty{1.22}{kg} \ce{CO2} per kilogram aluminium.
\textbf{25.} De minimumspanning zou met ongeveer \qty{1.03}{V} stijgen, maar de anode
zou zuurstof afgeven in plaats van koolstofdioxide en niet verbruikt worden.
\textbf{26.} De cel verbruikt $\boldsymbol{\approx \qty{13}{kW.h}}$ elektrische
energie per kilogram aluminium.
\end{solution}
